\chapter{Топологические пространства}
\section{Топология.}
\begin{definition}\label{def:1.1.1}
    Пусть $X$ -- некоторое непустое множество, $\tau \subset 2^X$ и справедливы утверждения A1 -- A3 про $\tau\colon$
    \begin{enumerate}
        \item[A1] $\emptyset, X \in \tau$.
        \item[A2] Пусть $I$ -- произвольное индексное множество.
        Тогда $\{U_\alpha\}_{\alpha \in I} \subset \tau \hookrightarrow
        \bigcup_{\alpha \in I} U_\alpha \in \tau$ (то есть $\tau$ замкнута относительно произвольных объединений).
        \item[A3] Пусть $n \in \N$. $\{U_i\}_{i = 1}^n \subset \tau \hookrightarrow \bigcap_{i = 1}^n U_i \in \tau$ (то есть $\tau$ замкнуто относительно конечных пересечений).
    \end{enumerate}
    Тогда пара $(X, \tau)$ называется топологическим пространством, $\tau$ -- топологией на $X$, а элементы $\tau$ называются открытыми ($\tau$-открытыми в случае если топология может пониматься неоднозначно) множествами.
\end{definition}
\begin{definition}
     Пусть $(X, \tau)$ -- топологическое пространство и $x \in A \in \tau$.
     Будем называть $A$ окрестностью точки $x$ (или просто окрестностью) и обозначать как $U(x), V(x), W(x)$.
\end{definition}
\begin{note}
    Поскольку $X \in \tau$ то всякая точка имеет хотя бы одну окрестность.
\end{note}
\begin{examples}
    На множестве $X$ всегда можно задать несколько топологий:
    \begin{enumerate}
        \item Топология $\tau_{\text{сильн.}} = 2^X$ называется сильной топологией. \footnote{Может называться дискретной топологией.}
        \item Топологию $\tau_{\text{слаб.}} =\{\emptyset, X\}$ будем называть слабой топологией. \footnote{Может называться антидискретной топологией.}
    \end{enumerate}
    Названия <<сильная>> и <<слабая>> топология обусловлены введением на множестве всех топологий частичного порядка по включению.
    \begin{definition}
        В частности, если $\tau_1$ и $\tau_2$ -- топологии на $X$, то будем говорить что $\tau_1$ слабее $\tau_2$ если $\tau_1 \subset \tau_2$.
        Аналогично, если $\tau_2 \supset \tau_1$, то $\tau_2$ -- сильнее $\tau_1$.
    \end{definition}
    В таких терминах $\tau_{\text{сильн.}}$ -- самая сильная топология, а $\tau_{\text{слаб.}}$ -- самая слабая топология.
\end{examples}
\begin{notet}
    В дальнейшем вместо того, чтобы писать <<$(X, \tau)$ -- топологическое пространство>> будет употребляться запись <<$(X, \tau) \in \cat{Top}$>>
    Под $\cat{Top}$ стоит понимать класс (категорию) всех топологических пространств.
\end{notet}
\begin{definition}
    Пусть $(X, \tau) \in \cat{Top}$.
    Рассмотрим семейство множеств $\phi \subset 2^X\colon$
    \[
        \phi = \{ A \subset X \mid X \setminus A \in \tau \}.
    \]
    Элементы $\phi$ называются (топологически) замкнутыми ($\tau$-замкнутыми, если может пониматься неоднозначно) множествами.
\end{definition}
\begin{note}
    Топологическое пространство можно определять не через задание семейства открытых множеств, а через задание семейства замкнутых множеств, то есть такого $\phi$, что
    \begin{enumerate}
        \item $\emptyset, X \in \phi$.
        \item $\{U_\alpha\}_{\alpha \in I} \subset \phi \hookrightarrow \bigcap_{\alpha \in I} U_\alpha \in \phi$.
        \item $\{U_i\}_{i = 1}^n \subset \phi \hookrightarrow \bigcup_{i = 1}^n U_i \in \phi$.
    \end{enumerate}
    Тогда семейство множеств, состоящее из дополнений множеств из $\phi$ будет являться топологией на $X$.
\end{note}
\begin{notet}
    Если у нас задано некоторое топологическое пространство $(X, \tau)$, то $\phi$ будет обозначать семейство его замкнутых множеств. \\
    Если у нас $(X, \tau) \in \cat{Top}$ (и нет других топологических пространств или к топологии $X$ мы явно не обращаемся) будем считать что $\tau$ -- топология на $X$ (без явного определения) и писать просто $X \in \cat{Top}$.
\end{notet}
\begin{definition}
    Если $(X, \tau) \in \cat{Top}$ и $S \subset X$, то будет называть (топологическим) замыканием (возможно, $\tau$-замыканием) множества $S$
    \[
        [S]_{\tau} = \bigcap \{A \mid S \subset A \in \phi\}.
    \]
     Заметим, что $[S]_\tau$ -- замкнуто, как пересечение произвольного числа замкнутых. \\
     Поскольку $X$ -- замкнуто, то $[S]_\tau$ -- корректно определено. \\
    Также, нетрудно заметить, что если $S \in \phi$, то $[S]_\tau = S$.
\end{definition}
\begin{definition}
    Пусть $(X, \tau) \in \cat{Top}$ и $S \subset X$.
    Тогда $x \in X$ называется точкой прикосновения множества $S$ если
    \[
        \forall U(x) \in \tau \hookrightarrow U(x) \cap S \neq \emptyset.
    \]
\end{definition}
\begin{lemma}
    Пусть $X \in \cat{Top}$ и $A \subset X$.
    Следующие утверждения эквивалентны:
    \begin{enumerate}
        \item $A \in \tau$.
        \item Всякая точка $A$ является внутренней, то есть входит в $A$ вместе с некоторой своей окрестностью.
        Формально:
        \[
            \forall x \in A \ \exists U(x) \in \tau \colon U(x) \subset A.
        \]
    \end{enumerate}
\end{lemma}
\begin{proof}
    Пусть $A$ открыто.
    Тогда для всякой точки $x \in A$ само множество $A$ является открытой окрестностью и всякая точка является внутренней. \\
    Пусть всякая точка $A$ -- внутренняя.
    Тогда
    \[
        A = \bigcup_{x \in A} \{x\} \subset \bigcup_{x \in A} U(x) \subset A,
    \]
    где $\forall x \in A$ окрестность $U(x)$ вложена в $A$. \\
    Но тогда $A$ -- объединение открытых множеств и, следовательно, открыто.
\end{proof}
\begin{proposition}
    Пусть $(X, \tau) \in \cat{Top}$ и $S \subset X$.
    Следующие условия эквивалентны:
    \begin{enumerate}
        \item $S \in \phi$.
        \item $S$ содержит все свои точки прикосновения.
    \end{enumerate}
    Более того, замыкание множества допускает следующее описание:
    \[
        [S]_\tau = \{x \in X \mid x \text{ -- точка прикосновения }S\}.
    \]
\end{proposition}
\begin{proof}
    Пусть $S \in \phi$ и $x \in X$ -- его точка прикосновения.
    Предположим, что $x \notin S$.
    Тогда $x \in X \setminus S \in \tau$ -- и возникает противоречие с определением точки прикосновения, так как $X \setminus S$ -- окрестность $x$, не пересекающаяся с $S$.
    Значит $S$ содержит все свои точки прикосновения.

    Пусть $S$ содержит все свои точки прикосновения.
    Рассмотрим точку $x \in X \setminus S$.
    Заметим, что она является внутренней для $X \setminus S$.
    И действительно: так как она не является точкой прикосновения $S$, то существует $U(x) \subset X \setminus S$.
    Тогда всякая точка $X \setminus S$ -- внутренняя и $X \setminus S$ открыто по предыдущей лемме, а значит $S$ является замкнутым по определению.

    Покажем справедливость следующего утверждения:
    \[
         [S]_\tau = \{x \in X \mid x \text{ -- точка прикосновения }S\}.
    \]
    Заметим, что $S \subset [S]_\tau$.
    Тогда всякая точка прикосновения $S$ является и точкой прикосновения $ [S]_\tau$.
    Поскольку $[S]_\tau$ -- замкнуто, то оно содержит все свои точки прикосновения и, в частности, содержит все точки прикосновения $S$.

    С другой стороны, пусть $x \in [S]_\tau$ не является точкой прикосновения $S$.
    Тогда $\exists U(x) \in \tau \colon U(x) \cap S = \emptyset$, а значит $U(x) \subset X \setminus S$ и $S \subset X \setminus U(x)$.
    Но $X \setminus U(x)$ -- замкнуто и содержит $S$, а следовательно и $[S]_\tau$ вложено в $X \setminus U(x)$.
    Приходим к противоречию: с одной стороны $x \in [S]_\tau \subset X \setminus U(x)$, а с другой $x \notin X \setminus U(x)$.
    Значит $x$ является точкой прикосновения $S$.
\end{proof}
\begin{definition}
    Пусть $X \in \cat{Top}$ и $\{x_n\}_{n = 1}^{\infty} \subset X$ -- последовательность элементов $X$.
    Будем говорить что $x_n \xrightarrow{\tau} x \in X$ если
    \[
        \forall U(x) \in \tau \ \exists N \in \N  \ \forall n \geq N \hookrightarrow x_n \in U(x).
    \]
\end{definition}
\begin{examples}
    Заметим, что в $\tau_{\text{слаб.}}$ (то есть в слабейшей топологии) все последовательности сходятся ко всем точкам. \\
    В сильнейшей топологии $\tau_{\text{сильн.}}$ сходятся только стационарные последовательности.
\end{examples}
\begin{definition}
    Пусть $X \in \cat{Top}$ и $S \subset X$.
    Точка $x \in X$ называется секвенциальной точкой прикосновения, если $\exists \{x_n\}_{n = 1}^{\infty} \subset S\colon$
    \[
        x_n \xrightarrow{\tau} x, n \rightarrow \infty.
    \]
\end{definition}
\begin{definition}
    Пусть $X \in \cat{Top}$ и $S \subset X$.
    \begin{enumerate}
        \item Будем говорить что $S$ -- секвенциально замкнуто, если оно содержит все свои секвенциальные точки прикосновения.
        \item Будем называть секвенциальным замыканием множества $S$ множество всех секвенциальных точек прикосновения множества $S$ и обозначать как $[S]_{seq}$.
    \end{enumerate}
\end{definition}
\begin{note}
    Заметим, что всякая секвенциальная точка прикосновения является и (топологической)\footnote{По умолчанию под замкнутыми понимаются именно топологически замкнутые множества} точкой прикосновения. \\
    Следовательно всякое (топологически) замкнутое множество является и секвенциально замкнутым. \\
    Более того, $\forall S \subset X \hookrightarrow [S]_{seq} \subset [S]_\tau$.
\end{note}
\begin{example}
    Рассмотрим $(\R, \tau)$\footnote{Роман Викторович называет эту топологию топологией Зарисского. Строго говоря, это неверно: в топологии Зарисского на $\mathbb{A}^n$ (где $\mathbb{A}^n$ -- аффинное пространство над, обычно, алгебраически замкнутым полем $K$) замкнутыми объявляются множества общих нулей множеств полиномов от $n$ переменных.
    В случае $\R$ нулями полиномов от одной переменной являются конечные множества и всё $\R$, поэтому таким образом определённая топология не совпадает с топологией Зарисского.
    }, где $\tau$ состоит из дополнений счётных множеств и пустого, то есть
    \[
        V \in \tau \Longleftrightarrow V = \emptyset \text{ или } |X \setminus V| \leq \aleph_0.
    \]
    Заметим, что в такой топологии $\tau$ сходятся только стационарные последовательности.
    И действительно: если $\{x_n\}_{n = 1}^{\infty} \subset \R$ и $x_n \xrightarrow{\tau} x \in \R$, то так как $W = \{x_n \mid n \in \N\} \setminus \{x\}$ -- не более чем счётно, то его дополнение $X \setminus W$ -- открыто и является окрестностью $x$.
    С другой стороны, если $x_n$ -- нестационарна, то в $X \setminus W$ не лежит бесконечное число элементов последовательности и, следовательно, последовательность не может сходиться к $x$.
    Значит $\{x_n\}$ -- стационарна, и, таким образом, всякое множество является секвенциально замкнутым (а секвенциальным замыканием всякого множества является оно же). \\
    С другой стороны, если мы рассмотрим произвольное континуальное множество $F$, то его топологическим замыканием будет $\R$.
    И действительно: все замкнутые в $\tau$ множества (кроме $\R$) не более чем счётны и, следовательно, не могут содержать $F$.
    Значит единственным замкнутым множеством, содержащим $F$, является $\R$ и $[F]_\tau = \R$ по определению.
\end{example}
\begin{note}
    Существуют такие пространства $X \in \cat{Top}$, в которых секвенциальные замыкания множеств могут не быть секвенциально замкнутыми, то есть $\left[[S]_{seq}\right]_{seq} \supsetneq [S]_{seq}$. \\
    В частности, в $\R^{[0, 1]}$ (пространство функций на отрезке $[0, 1]$) с его тихоновской топологией множество $C[0, 1]$ (непрерывных функций на отрезке $[0, 1]$ с индуцированной топологией) не является секвенциально замкнутым, а всякое его секвенциальное замыкание строго вложено в следующее.
\end{note}
\begin{note}
    Нетрудно заметить, что семейство секвенциально-замкнутых множество $\phi_{seq}$ для некоторого топологического пространства $X$ является семейством замкнутых множеств:
    \begin{enumerate}
        \item Очевидно, что и $\emptyset$, и $X$ -- секвенциально замкнуты.
        \item Если $\{U_\alpha\}_{\alpha \in I} \subset \phi_{seq}$, то $V = \bigcap_{\alpha \in I} U_\alpha \in \phi_{seq}$, поскольку если $x$ -- секвенциальная точка прикосновения для $V$, то, в частности, она является секвенциальной точкой прикосновения для $U_\alpha$ при всяком $\alpha$ -- а следовательно лежит в $U_\alpha$.
        Но тогда она лежит и в $V$ -- а следовательно $V$ секвенциально замкнуто.
        \item Если $\{U_i\}_{i = 1}^n \subset \phi_{seq}$, то $W = \bigcup_{i = 1}^n U_i \in \phi_{seq}$, поскольку если $x \in X$ -- секвенциальная точка прикосновения, то существует $\{x_n\} \subset W\colon x_n \xrightarrow{\tau} x$.
        Но поскольку $W$ -- объединение конечного числа множеств $U_i$, то существует подпоследовательность $\{x_{n_k}\}$, которая целиком лежит в некотором $U_s$ -- а следовательно и в $U_s$, и в $W$.
    \end{enumerate}
    Таким образом семейство дополнений множеств из $\phi_{seq}$ является топологией на $X$, которая не слабее $\tau$. \\
    В частности, для предыдущего примера (то есть $(\R, \tau)$) справедливо, что $\tau_{seq} = 2^{\R}$.
\end{note}
